\section{Experimental Setup}
\label{sec:setup}

\paragraph{Datasets and groups.}
We use \V{nDatasets} image classification datasets (\cref{tab:datasets}).
The \V{nOther} FGIR datasets cover birds, aircraft, cars, dogs, flowers, food, pets, vegetables and fruit, anime characters, and natural species~\cite{wah2011cub,vanhorn2015nabirds,maji2013fine,krause2013cars,khosla2011dogs,nilsback2008flowers,bossard2014food,parkhi2012pets,hou2017vegfru,vanhorn2018inat} \TBD{citation for Moe}.
The other \V{nSc} are the four soybean cultivar datasets and the cotton cultivar dataset of the ultra-fine-grained benchmark~\cite{yu2021benchmark}.
We call them the \SC{} group and the remaining \V{nOther} datasets the other group.
The grouping follows data source, not a measured property.
The two groups also differ in baseline accuracy: averaged over datasets, baseline top-1 is \V{fpAccSc}\% for the \SC{} group and \V{fpAccOther}\% for the other group (\cref{tab:datasets}).

\begin{table}[t]
\centering
\caption{Datasets, their group, split sizes and number of classes $C$. FP is the baseline top-1, the mean over the nine model and regime settings at the AMP FP16 baseline.}
\label{tab:datasets}
\setlength{\tabcolsep}{2.5pt}
\small
\begin{tabular}{llrrrr}
\toprule
Dataset & Group & Train $N$ & Test $N$ & $C$ & FP (\%) \\
\midrule
Aircraft & Other & \dummy{1{,}000} & 3{,}333 & 100 & 78.4 \\
Cars & Other & \dummy{1{,}037} & 8{,}041 & 196 & 82.7 \\
CUB & Other & \dummy{1{,}074} & 5{,}794 & 200 & 89.1 \\
Dogs & Other & \dummy{1{,}111} & 8{,}580 & 120 & 88.5 \\
Flowers & Other & \dummy{1{,}148} & 6{,}149 & 102 & 98.8 \\
Food & Other & \dummy{1{,}185} & 25{,}250 & 101 & 90.4 \\
iNat17 & Other & \dummy{1{,}222} & 95{,}986 & 5{,}089 & 68.7 \\
Moe & Other & \dummy{1{,}259} & 2{,}943 & 173 & 91.4 \\
NABirds & Other & \dummy{1{,}296} & 24{,}633 & 555 & 86.7 \\
Pets & Other & \dummy{1{,}333} & 3{,}669 & 37 & 92.3 \\
VegFru & Other & \dummy{1{,}370} & 116{,}931 & 292 & 93.2 \\
SoyAgeing & Soy+cotton & \dummy{1{,}407} & 4{,}950 & 198 & 44.2 \\
SoyGene & Soy+cotton & \dummy{1{,}444} & 11{,}143 & 1{,}110 & 37.8 \\
SoyGlobal & Soy+cotton & \dummy{1{,}481} & 5{,}814 & 1{,}938 & 20.4 \\
SoyLocal & Soy+cotton & \dummy{1{,}518} & 600 & 200 & 27.1 \\
Cotton & Soy+cotton & \dummy{1{,}555} & 240 & 80 & 32.8 \\
\bottomrule
\end{tabular}
\end{table}


\paragraph{Backbones and checkpoint regimes.}
We use ViT-B/16~\cite{dosovitskiy2021vit}, Swin-B~\cite{liu2021swin} and BEiTv2-B~\cite{peng2022beitv2} at $224\times224$ input, each trained on every dataset under three regimes.
FT fine-tunes the whole network, FZ trains only the classifier on a frozen backbone, and CAL fine-tunes with counterfactual attention learning~\cite{rao2021cal} \TBD{confirm regime definitions, pretrained weights and training hyperparameters from the run records}.
CAL predictions aggregate cropped and flipped views, while its features come from the single test image.
Each dataset, backbone and regime defines one context, giving 144 contexts with \V{testRows} test images per backbone and regime.
Test images are resized to 300 pixels and center-cropped to 224.

\paragraph{Quantization.}
We quantize the weights of every linear layer in the backbone with HQQ and keep the classifier head at full precision \TBD{HQQ version, and the number of quantized layers per backbone}.
Activations are not quantized.
The grid covers bit widths $b \in \{1, 1.58, 2, 3, 4, 5, 6, 8\}$ and group sizes $g \in \{8, 32, 128, 512\}$, giving \V{nConditions} conditions; we write W3/G128 for $b=3$ and $g=128$.
The baseline of each context is the same checkpoint evaluated without HQQ under automatic mixed precision with FP16, on the same test images, with one inference seed.
The grid holds \V{cells} quantized results over \V{contexts} contexts.
VegFru has no Swin-B results and only 10 of 32 ViT-B/16 CAL results, so grid analyses use the \V{gridDatasets} complete datasets.

\paragraph{Accuracy retention.}
For a context with baseline top-1 accuracy $a_{\mathrm{FP}}$ and quantized accuracy $a_{\mathrm{Q}}$, the accuracy retention is
\begin{equation}
  r = a_{\mathrm{Q}} / a_{\mathrm{FP}},
  \label{eq:retention}
\end{equation}
and the relative loss is $1-r$.
Group retention is the mean of $r$ over the datasets of a group, each weighted equally.
The retention gap is the other group's retention minus the \SC{} group's.

\paragraph{Dataset difficulty.}
We fix the difficulty measure before looking at its correlates.
The primary conditions are W3 and W4 at group sizes 8 and 128, because retention saturates near zero at 1 and 1.58 bits and near one at 6 and 8 bits.
Within each backbone and regime, we rank the datasets by relative loss in each primary condition as percentiles, and average the four percentiles into a difficulty score.
Kendall's $W$ measures how well the four conditions agree on the ranking; we treat the score as a common difficulty only when $W \ge 0.6$.

\paragraph{Decision margin.}
For logits $\bm{z}$ of an image $x$ under the baseline, the normalized decision margin is
\begin{equation}
  m(x) = \frac{z_{(1)} - z_{(2)}}{\sigma(\bm{z})},
  \label{eq:margin}
\end{equation}
where $z_{(1)}$ and $z_{(2)}$ are the two largest logits and $\sigma(\bm{z})$ is their standard deviation over classes.
The normalization makes margins comparable across checkpoints with different logit scales.
The margin measures the distance to the second-ranked class only, while quantization can also make a lower-ranked class the winner.
The vulnerable-tail margin $q_{10}$ of a context is the 10th percentile of $m(x)$ over its test images.
We prespecify $q_{10}$ as the primary margin predictor.
It gives a difficulty signal when its Spearman correlation with the difficulty score is $\rho \le -0.5$ with a one-sided dataset permutation $p < 0.05$.
The correlation must also be negative in at least three of the four primary conditions and after leaving out any one dataset.
Baseline accuracy and the logarithm of the number of classes are the prespecified controls.
We test whether adding $q_{10}$ to them lowers the leave-one-dataset-out error of a linear model of the difficulty score.

\paragraph{Paired W3/G128 analysis.}
At W3/G128 we export the features after every block and the final logits of both models for all 144 contexts.
A swap is an image whose predicted class changes after quantization.
Damage is a swap from correct to wrong and rescue a swap from wrong to correct; both are rates over all test images.
Layer-wise similarity is the linear CKA~\cite{kornblith2019similarity} between the centered baseline and quantized features $X$ and $Y$,
\begin{equation}
  \mathrm{CKA}(X, Y) = \frac{\lVert X^{\top} Y \rVert_F^2}{\lVert X^{\top} X \rVert_F \, \lVert Y^{\top} Y \rVert_F}.
  \label{eq:cka}
\end{equation}
On $\ell_2$-normalized final features, $D_{\mathrm{intra}}$ is the mean cosine distance between images of one class and $D_{\mathrm{inter}}$ between images of two classes, each averaged with classes weighted equally.
Separability is $D_{\mathrm{inter}} - D_{\mathrm{intra}}$.
The H-score is $\mathrm{tr}\big(S_b (S_w + \lambda I)^{-1}\big)$ with between-class and within-class scatter $S_b$ and $S_w$ and $\lambda = 0.01\,\mathrm{tr}(S_w)/D$ for feature dimension $D$.

\paragraph{Margin matching.}
To ask how much of the retention gap follows from the different margin distributions of the groups, we reweight the test images of both groups to a common margin distribution.
We bin $m(x)$ with width 0.05, weight the datasets within each group equally, and take the bin-wise minimum of the two group histograms as the common target.
Each image receives the ratio of target to group mass in its bin, and group retention becomes weighted quantized accuracy over weighted baseline accuracy.
The weights depend only on the baseline margin.
We report the effective sample size (ESS) and the overlap of the two histograms beside the matched gap.

\paragraph{Logit offsets and competing classes.}
Quantization may shift the logits of some classes more than others.
We estimate the mean row-centered logit change of each class on four folds of a stratified five-fold split and subtract it on the held-out fold.
Five controls apply the same correction with the class assignment of the offsets shuffled.
This correction uses test logits and is a diagnostic, not a deployable method.
To ask whether competing classes predict damage beyond the margin, we count the classes whose baseline logit lies within 0.25, 0.5 and 1.0 logit standard deviations of the winner.
Two logistic models predict damage among baseline-correct images: one from a nonlinear basis of $m(x)$ and the log number of classes, and one that adds the counts.
Both use L2 regularization with $C=100$ and are evaluated leave-one-dataset-out within each backbone and regime, by Brier score.
